\chapter{End-To-End Terminal Scenarios}\label{record-015}

\noindent{\footnotesize\color{BookMuted}\path{docs/guides/terminal-scenarios.md}\par}\smallskip

This guide walks through Bijou from small one-shot terminal output to full-screen TUI apps and shader-driven visual renderers. It is scenario-first: pick the thing you want to build, then follow the package boundary and control flow that fits the job.

Use this guide when you are asking one of these questions:

\begin{itemize}
\tightlist
\item
  I want to print some colors.
\item
  I want to make an interactive form.
\item
  I want to make a basic TUI app.
\item
  I want to make a Braille ray trace renderer.
\end{itemize}

For a live tour of the repository product surface, run:

\begin{CodeBlock}
npm run dogfood
\end{CodeBlock}

For focused component and preview work, run:

\begin{CodeBlock}
npm run blocklab
\end{CodeBlock}

\section{Package Stack}\label{record-015-package-stack}

Bijou is intentionally layered. Core rendering and contracts are portable; interactive runtime behavior lives in the TUI package; host I/O lives behind Node adapters.

\begin{center}
\resizebox{\linewidth}{!}{\begin{tikzpicture}[box/.style={draw=BookHeading,fill=BookPanel,rounded corners=1pt,text width=38mm,inner sep=5pt,align=left,font=\footnotesize},edge/.style={-{Latex[length=2mm]},draw=BookHeading}]
\node[box] (App) at (0.0,0.0) {Your app code};
\node[box] (Core) at (-5.0,-2.6) {@flyingrobots/bijou / components, surfaces, forms, blocks, themes};
\node[box] (Tui) at (0.0,-2.6) {@flyingrobots/bijou-tui / TEA runtime, input, layout, canvas};
\node[box] (Node) at (5.0,-2.6) {@flyingrobots/bijou-node / terminal IO, style adapter, host runner};
\node[box] (Terminal) at (2.5,-5.2) {Terminal stdin/stdout};
\node[box] (Surface) at (-2.5,-5.2) {Surface or string output};
\draw[edge,] (App) to (Core);
\draw[edge,] (App) to (Tui);
\draw[edge,] (App) to (Node);
\draw[edge,bend left=15] (Tui) to (Core);
\draw[edge,bend left=15] (Node) to (Core);
\draw[edge,] (Node) to (Terminal);
\draw[edge,] (Core) to (Surface);
\draw[edge,] (Tui) to (Surface);
\draw[edge,] (Surface) to (Node);
\end{tikzpicture}}
\end{center}

\begin{center}
\resizebox{\linewidth}{!}{\begin{tikzpicture}[box/.style={draw=BookHeading,fill=BookPanel,rounded corners=1pt,text width=38mm,inner sep=5pt,align=left,font=\footnotesize},edge/.style={-{Latex[length=2mm]},draw=BookHeading}]
\node[box] (BijouContext) at (-5.0,0) {\textbf{Bijou\allowbreak{}Context}\\\rule{\linewidth}{.3pt}\\+runtime\\+io\\+style\\+theme\\+semantic(name)\\+status(name)\\+border(name)};
\node[box] (Surface) at (0.0,0) {\textbf{Surface}\\\rule{\linewidth}{.3pt}\\+width\\+height\\+get(x, y)\\+set(x, y, cell)};
\node[box] (App) at (5.0,0) {\textbf{App}\\\rule{\linewidth}{.3pt}\\+init()\\+update(msg, model)\\+view(model)};
\node[box] (Cmd) at (-2.5,-4) {\textbf{Cmd}\\\rule{\linewidth}{.3pt}\\+execute(emit, caps)};
\node[box] (NodeHost) at (2.5,-4) {\textbf{Node\allowbreak{}Host}\\\rule{\linewidth}{.3pt}\\+start\allowbreak{}App(app, options)\\+create\allowbreak{}Node\allowbreak{}Context()};
\draw[edge,bend left=15] (BijouContext) to (Surface);
\draw[edge,bend left=15] (App) to (Surface);
\draw[edge,] (App) to (Cmd);
\draw[edge,] (NodeHost) to (App);
\draw[edge,] (NodeHost) to (BijouContext);
\end{tikzpicture}}
\end{center}

\section{Choose A Scenario}\label{record-015-choose-a-scenario}

\begin{center}
\resizebox{\linewidth}{!}{\begin{tikzpicture}[box/.style={draw=BookHeading,fill=BookPanel,rounded corners=1pt,text width=38mm,inner sep=5pt,align=left,font=\footnotesize},edge/.style={-{Latex[length=2mm]},draw=BookHeading}]
\node[box] (Start) at (0.0,0.0) {What are you building?};
\node[box] (Colors) at (-7.5,-2.6) {One-shot colored CLI output};
\node[box] (Form) at (-2.5,-2.6) {Prompt or form flow};
\node[box] (TUI) at (2.5,-2.6) {Full-screen TUI};
\node[box] (Visual) at (7.5,-2.6) {Procedural visual renderer};
\node[box] (Core) at (-2.5,-5.2) {@flyingrobots/bijou + @flyingrobots/bijou-node};
\node[box] (Runtime) at (2.5,-5.2) {@flyingrobots/bijou-tui + @flyingrobots/bijou-node};
\draw[edge,] (Start) to (Colors);
\draw[edge,] (Start) to (Form);
\draw[edge,] (Start) to (TUI);
\draw[edge,] (Start) to (Visual);
\draw[edge,] (Colors) to (Core);
\draw[edge,] (Form) to (Core);
\draw[edge,] (TUI) to (Runtime);
\draw[edge,] (Visual) to (Runtime);
\end{tikzpicture}}
\end{center}

{\def\LTcaptype{none} % do not increment counter
\begin{longtable}[]{@{}
  >{\raggedright\arraybackslash}p{(\linewidth - 4\tabcolsep) * \real{0.3333}}
  >{\raggedright\arraybackslash}p{(\linewidth - 4\tabcolsep) * \real{0.3333}}
  >{\raggedright\arraybackslash}p{(\linewidth - 4\tabcolsep) * \real{0.3333}}@{}}
\toprule\noalign{}
{S\allowbreak{}c\allowbreak{}e\allowbreak{}n\allowbreak{}a\allowbreak{}r\allowbreak{}i\allowbreak{}o}
 & {P\allowbreak{}r\allowbreak{}i\allowbreak{}m\allowbreak{}a\allowbreak{}r\allowbreak{}y} {p\allowbreak{}a\allowbreak{}c\allowbreak{}k\allowbreak{}a\allowbreak{}g\allowbreak{}e\allowbreak{}s}
 & {M\allowbreak{}a\allowbreak{}i\allowbreak{}n} {o\allowbreak{}u\allowbreak{}t\allowbreak{}p\allowbreak{}u\allowbreak{}t}
 \\
\midrule\noalign{}
\endhead
\bottomrule\noalign{}
\endlastfoot
{P\allowbreak{}r\allowbreak{}i\allowbreak{}n\allowbreak{}t} {c\allowbreak{}o\allowbreak{}l\allowbreak{}o\allowbreak{}r\allowbreak{}s} & \texttt{@\allowbreak{}f\allowbreak{}l\allowbreak{}y\allowbreak{}i\allowbreak{}n\allowbreak{}g\allowbreak{}r\allowbreak{}o\allowbreak{}b\allowbreak{}o\allowbreak{}t\allowbreak{}s\allowbreak{}/\allowbreak{}b\allowbreak{}i\allowbreak{}j\allowbreak{}o\allowbreak{}u}{,} \texttt{@\allowbreak{}f\allowbreak{}l\allowbreak{}y\allowbreak{}i\allowbreak{}n\allowbreak{}g\allowbreak{}r\allowbreak{}o\allowbreak{}b\allowbreak{}o\allowbreak{}t\allowbreak{}s\allowbreak{}/\allowbreak{}b\allowbreak{}i\allowbreak{}j\allowbreak{}o\allowbreak{}u\allowbreak{}-\allowbreak{}n\allowbreak{}o\allowbreak{}d\allowbreak{}e} & {S\allowbreak{}t\allowbreak{}r\allowbreak{}i\allowbreak{}n\allowbreak{}g\allowbreak{}s} {w\allowbreak{}i\allowbreak{}t\allowbreak{}h} {h\allowbreak{}o\allowbreak{}s\allowbreak{}t} {s\allowbreak{}t\allowbreak{}y\allowbreak{}l\allowbreak{}i\allowbreak{}n\allowbreak{}g} \\
{I\allowbreak{}n\allowbreak{}t\allowbreak{}e\allowbreak{}r\allowbreak{}a\allowbreak{}c\allowbreak{}t\allowbreak{}i\allowbreak{}v\allowbreak{}e} {f\allowbreak{}o\allowbreak{}r\allowbreak{}m} & \texttt{@\allowbreak{}f\allowbreak{}l\allowbreak{}y\allowbreak{}i\allowbreak{}n\allowbreak{}g\allowbreak{}r\allowbreak{}o\allowbreak{}b\allowbreak{}o\allowbreak{}t\allowbreak{}s\allowbreak{}/\allowbreak{}b\allowbreak{}i\allowbreak{}j\allowbreak{}o\allowbreak{}u}{,} \texttt{@\allowbreak{}f\allowbreak{}l\allowbreak{}y\allowbreak{}i\allowbreak{}n\allowbreak{}g\allowbreak{}r\allowbreak{}o\allowbreak{}b\allowbreak{}o\allowbreak{}t\allowbreak{}s\allowbreak{}/\allowbreak{}b\allowbreak{}i\allowbreak{}j\allowbreak{}o\allowbreak{}u\allowbreak{}-\allowbreak{}n\allowbreak{}o\allowbreak{}d\allowbreak{}e} & {P\allowbreak{}r\allowbreak{}o\allowbreak{}m\allowbreak{}p\allowbreak{}t} {r\allowbreak{}e\allowbreak{}s\allowbreak{}u\allowbreak{}l\allowbreak{}t\allowbreak{}s} \\
{B\allowbreak{}a\allowbreak{}s\allowbreak{}i\allowbreak{}c} {T\allowbreak{}U\allowbreak{}I} {a\allowbreak{}p\allowbreak{}p} & \texttt{@\allowbreak{}f\allowbreak{}l\allowbreak{}y\allowbreak{}i\allowbreak{}n\allowbreak{}g\allowbreak{}r\allowbreak{}o\allowbreak{}b\allowbreak{}o\allowbreak{}t\allowbreak{}s\allowbreak{}/\allowbreak{}b\allowbreak{}i\allowbreak{}j\allowbreak{}o\allowbreak{}u\allowbreak{}-\allowbreak{}t\allowbreak{}u\allowbreak{}i}{,} \texttt{@\allowbreak{}f\allowbreak{}l\allowbreak{}y\allowbreak{}i\allowbreak{}n\allowbreak{}g\allowbreak{}r\allowbreak{}o\allowbreak{}b\allowbreak{}o\allowbreak{}t\allowbreak{}s\allowbreak{}/\allowbreak{}b\allowbreak{}i\allowbreak{}j\allowbreak{}o\allowbreak{}u\allowbreak{}-\allowbreak{}n\allowbreak{}o\allowbreak{}d\allowbreak{}e} & \texttt{S\allowbreak{}u\allowbreak{}r\allowbreak{}f\allowbreak{}a\allowbreak{}c\allowbreak{}e} {f\allowbreak{}r\allowbreak{}a\allowbreak{}m\allowbreak{}e\allowbreak{}s} \\
{B\allowbreak{}r\allowbreak{}a\allowbreak{}i\allowbreak{}l\allowbreak{}l\allowbreak{}e} {r\allowbreak{}a\allowbreak{}y} {t\allowbreak{}r\allowbreak{}a\allowbreak{}c\allowbreak{}e} {r\allowbreak{}e\allowbreak{}n\allowbreak{}d\allowbreak{}e\allowbreak{}r\allowbreak{}e\allowbreak{}r} & \texttt{@\allowbreak{}f\allowbreak{}l\allowbreak{}y\allowbreak{}i\allowbreak{}n\allowbreak{}g\allowbreak{}r\allowbreak{}o\allowbreak{}b\allowbreak{}o\allowbreak{}t\allowbreak{}s\allowbreak{}/\allowbreak{}b\allowbreak{}i\allowbreak{}j\allowbreak{}o\allowbreak{}u\allowbreak{}-\allowbreak{}t\allowbreak{}u\allowbreak{}i}{,} \texttt{@\allowbreak{}f\allowbreak{}l\allowbreak{}y\allowbreak{}i\allowbreak{}n\allowbreak{}g\allowbreak{}r\allowbreak{}o\allowbreak{}b\allowbreak{}o\allowbreak{}t\allowbreak{}s\allowbreak{}/\allowbreak{}b\allowbreak{}i\allowbreak{}j\allowbreak{}o\allowbreak{}u\allowbreak{}-\allowbreak{}n\allowbreak{}o\allowbreak{}d\allowbreak{}e} & {S\allowbreak{}h\allowbreak{}a\allowbreak{}d\allowbreak{}e\allowbreak{}r\allowbreak{}-\allowbreak{}r\allowbreak{}e\allowbreak{}n\allowbreak{}d\allowbreak{}e\allowbreak{}r\allowbreak{}e\allowbreak{}d} \texttt{S\allowbreak{}u\allowbreak{}r\allowbreak{}f\allowbreak{}a\allowbreak{}c\allowbreak{}e} \\
\end{longtable}
}

\section{Scenario 1: Print Some Colors}\label{record-015-scenario-1-print-some-colors}

Use core components and a Node context when the program writes a result and exits. This is the right lane for CLIs, scripts, release notes, status output, and other non-full-screen tools.

\begin{center}
\resizebox{\linewidth}{!}{\begin{tikzpicture}[actor/.style={draw=BookHeading,fill=BookPanel,text width=27mm,align=center,font=\small,inner sep=5pt}]
\node[actor] (Script) at (0.0,0) {Script};
\draw[dashed,BookMuted] (Script.south) -- (0.0,-7.3999999999999995);
\node[actor] (Node) at (4.5,0) {Node adapter};
\draw[dashed,BookMuted] (Node.south) -- (4.5,-7.3999999999999995);
\node[actor] (Core) at (9.0,0) {Bijou core};
\draw[dashed,BookMuted] (Core.south) -- (9.0,-7.3999999999999995);
\node[actor] (T) at (13.5,0) {Terminal};
\draw[dashed,BookMuted] (T.south) -- (13.5,-7.3999999999999995);
\draw[-{Latex[length=2mm]},BookHeading] (0.0,-1.7) -- node[above,fill=white,align=center,text width=4.05cm,font=\footnotesize] {create\allowbreak{}Node\allowbreak{}Context()} (4.5,-1.7);
\draw[-{Latex[length=2mm]},BookHeading] (0.0,-3.4) -- node[above,fill=white,align=center,text width=8.1cm,font=\footnotesize] {separator(), progress\allowbreak{}Bar(), gradient\allowbreak{}Text()} (9.0,-3.4);
\draw[dashed,-{Latex[length=2mm]},BookHeading] (9.0,-5.1) -- node[above,fill=white,align=center,text width=8.1cm,font=\footnotesize] {styled strings} (0.0,-5.1);
\draw[-{Latex[length=2mm]},BookHeading] (0.0,-6.8) -- node[above,fill=white,align=center,text width=12.15cm,font=\footnotesize] {console.log(...)} (13.5,-6.8);
\end{tikzpicture}}\end{center}

\begin{center}
\resizebox{\linewidth}{!}{\begin{tikzpicture}[box/.style={draw=BookHeading,fill=BookPanel,rounded corners=1pt,text width=38mm,inner sep=5pt,align=left,font=\footnotesize},edge/.style={-{Latex[length=2mm]},draw=BookHeading}]
\node[box] (SCRIPT) at (0.0,0.0) {SCRIPT};
\node[box] (BIJOU_CONTEXT) at (0.0,-2.6) {BIJOU\_CONTEXT};
\node[box] (STYLE_PORT) at (-2.5,-5.2) {STYLE\_PORT};
\node[box] (THEME) at (2.5,-5.2) {THEME};
\node[box] (TOKEN) at (2.5,-7.800000000000001) {TOKEN};
\node[box] (ANSI_STRING) at (-2.5,-7.800000000000001) {ANSI\_STRING};
\draw[edge,] (SCRIPT) to node[midway,fill=white,font=\scriptsize,text=BookMuted,align=center,text width=27mm] {creates / one : one} (BIJOU_CONTEXT);
\draw[edge,] (BIJOU_CONTEXT) to node[midway,fill=white,font=\scriptsize,text=BookMuted,align=center,text width=27mm] {owns / one : one} (STYLE_PORT);
\draw[edge,] (BIJOU_CONTEXT) to node[midway,fill=white,font=\scriptsize,text=BookMuted,align=center,text width=27mm] {owns / one : one} (THEME);
\draw[edge,] (THEME) to node[midway,fill=white,font=\scriptsize,text=BookMuted,align=center,text width=27mm] {resolves / one : zero or more} (TOKEN);
\draw[edge,] (STYLE_PORT) to node[midway,fill=white,font=\scriptsize,text=BookMuted,align=center,text width=27mm] {emits / one : zero or more} (ANSI_STRING);
\end{tikzpicture}}
\end{center}

\subsection{Code}\label{record-015-code}

\begin{CodeBlock}
import {
  alert,
  box,
  gradientText,
  progressBar,
  separator,
} from '@flyingrobots/bijou';
import { createNodeContext } from '@flyingrobots/bijou-node';

const ctx = createNodeContext();
const brandGradient = ctx.theme.theme.gradient.brand;

console.log(separator({ label: 'deploy', width: 64, ctx }));
console.log(ctx.style.styled(ctx.status('success'), 'OK'));
console.log(gradientText('Bijou terminal output', brandGradient, {
  style: ctx.style,
  noColor: ctx.theme.noColor,
}));
console.log(progressBar(72, { width: 32, showPercent: true, ctx }));
console.log(alert('Production deploy is ready.', { variant: 'success', ctx }));
console.log(box('Use core components for one-shot terminal output.', { ctx }));
\end{CodeBlock}

\subsection{Rule Of Thumb}\label{record-015-rule-of-thumb}

Use string-returning components for command output. Do not start the TUI runtime unless you need retained state, input routing, animation, or full-screen redraw.

\begin{center}
\resizebox{\linewidth}{!}{\begin{tikzpicture}[box/.style={draw=BookHeading,fill=BookPanel,rounded corners=1pt,text width=38mm,inner sep=5pt,align=left,font=\footnotesize},edge/.style={-{Latex[length=2mm]},draw=BookHeading}]
\node[box] (Need) at (0.0,0.0) {Need terminal output};
\node[box] (Input) at (0.0,-2.6) {Need ongoing input?};
\node[box] (Core) at (-2.5,-5.2) {Use core components and Node context};
\node[box] (Fullscreen) at (2.5,-5.2) {Need full-screen retained UI?};
\node[box] (Forms) at (-2.5,-7.800000000000001) {Use prompt/form helpers};
\node[box] (TUI) at (2.5,-7.800000000000001) {Use App + start\allowbreak{}App};
\draw[edge,] (Need) to (Input);
\draw[edge,] (Input) to node[midway,fill=white,font=\scriptsize,text=BookMuted,align=center,text width=27mm] {no} (Core);
\draw[edge,] (Input) to node[midway,fill=white,font=\scriptsize,text=BookMuted,align=center,text width=27mm] {yes} (Fullscreen);
\draw[edge,] (Fullscreen) to node[midway,fill=white,font=\scriptsize,text=BookMuted,align=center,text width=27mm] {no} (Forms);
\draw[edge,] (Fullscreen) to node[midway,fill=white,font=\scriptsize,text=BookMuted,align=center,text width=27mm] {yes} (TUI);
\end{tikzpicture}}
\end{center}

\section{Scenario 2: Make An Interactive Form}\label{record-015-scenario-2-make-an-interactive-form}

Use forms when the program should ask a bounded sequence of questions and then return to normal terminal output. Forms write through the context I/O port and return structured values.

\begin{center}
\resizebox{\linewidth}{!}{\begin{tikzpicture}[actor/.style={draw=BookHeading,fill=BookPanel,text width=27mm,align=center,font=\small,inner sep=5pt}]
\node[actor] (User) at (0.0,0) {User};
\draw[dashed,BookMuted] (User.south) -- (0.0,-12.5);
\node[actor] (Form) at (4.5,0) {Form};
\draw[dashed,BookMuted] (Form.south) -- (4.5,-12.5);
\node[actor] (Core) at (9.0,0) {Bijou forms};
\draw[dashed,BookMuted] (Core.south) -- (9.0,-12.5);
\node[actor] (Host) at (13.5,0) {Node context};
\draw[dashed,BookMuted] (Host.south) -- (13.5,-12.5);
\draw[-{Latex[length=2mm]},BookHeading] (0.0,-1.7) -- node[above,fill=white,align=center,text width=4.05cm,font=\footnotesize] {run setup flow} (4.5,-1.7);
\draw[-{Latex[length=2mm]},BookHeading] (4.5,-3.4) -- node[above,fill=white,align=center,text width=4.05cm,font=\footnotesize] {group(fields)} (9.0,-3.4);
\draw[-{Latex[length=2mm]},BookHeading] (9.0,-5.1) -- node[above,fill=white,align=center,text width=4.05cm,font=\footnotesize] {write prompts} (13.5,-5.1);
\draw[dashed,-{Latex[length=2mm]},BookHeading] (0.0,-6.8) -- node[above,fill=white,align=center,text width=12.15cm,font=\footnotesize] {keyboard input} (13.5,-6.8);
\draw[dashed,-{Latex[length=2mm]},BookHeading] (13.5,-8.5) -- node[above,fill=white,align=center,text width=4.05cm,font=\footnotesize] {answers} (9.0,-8.5);
\draw[dashed,-{Latex[length=2mm]},BookHeading] (9.0,-10.2) -- node[above,fill=white,align=center,text width=4.05cm,font=\footnotesize] {\{ cancelled, values \}} (4.5,-10.2);
\draw[-{Latex[length=2mm]},BookHeading] (4.5,-11.9) -- node[above,fill=white,align=center,text width=8.1cm,font=\footnotesize] {print summary} (13.5,-11.9);
\end{tikzpicture}}\end{center}

\begin{center}
\resizebox{\linewidth}{!}{\begin{tikzpicture}[box/.style={draw=BookHeading,fill=BookPanel,rounded corners=1pt,text width=38mm,inner sep=5pt,align=left,font=\footnotesize},edge/.style={-{Latex[length=2mm]},draw=BookHeading}]
\node[box] (GroupResult) at (-5.0,0) {\textbf{Group\allowbreak{}Result}\\\rule{\linewidth}{.3pt}\\+cancelled\\+values};
\node[box] (Field) at (0.0,0) {\textbf{Field}\\\rule{\linewidth}{.3pt}\\+title\\+placeholder\\+required\\+ctx};
\node[box] (FormHelpers) at (5.0,0) {\textbf{Form\allowbreak{}Helpers}\\\rule{\linewidth}{.3pt}\\+input(options)\\+select(options)\\+multiselect(options)\\+confirm(options)\\+group(fields)};
\draw[edge,bend left=15] (FormHelpers) to (Field);
\draw[edge,bend left=15] (FormHelpers) to (GroupResult);
\end{tikzpicture}}
\end{center}

\subsection{Code}\label{record-015-code-1}

\begin{CodeBlock}
import { box, confirm, group, input, multiselect, select } from '@flyingrobots/bijou';
import { initDefaultContext } from '@flyingrobots/bijou-node';

const ctx = initDefaultContext();

const result = await group({
  name: () => input({
    title: 'Project name:',
    placeholder: 'my-app',
    required: true,
    ctx,
  }),
  framework: () => select({
    title: 'Framework:',
    options: [
      { label: 'Express', value: 'express' },
      { label: 'Fastify', value: 'fastify' },
      { label: 'Hono', value: 'hono' },
      { label: 'None', value: 'none' },
    ],
    ctx,
  }),
  features: () => multiselect({
    title: 'Features:',
    options: [
      { label: 'TypeScript', value: 'typescript' },
      { label: 'ESLint', value: 'eslint' },
      { label: 'Docker', value: 'docker' },
    ],
    ctx,
  }),
  deploy: () => confirm({
    title: 'Set up deployment?',
    defaultValue: true,
    ctx,
  }),
});

if (result.cancelled) {
  console.log('Setup cancelled.');
} else {
  console.log(box([
    `Name: ${result.values.name}`,
    `Framework: ${result.values.framework}`,
    `Features: ${result.values.features.join(', ') || 'none'}`,
    `Deploy: ${result.values.deploy ? 'yes' : 'no'}`,
  ].join('\n'), { ctx }));
}
\end{CodeBlock}

\subsection{Data Shape}\label{record-015-data-shape}

\begin{center}
\resizebox{\linewidth}{!}{\begin{tikzpicture}[box/.style={draw=BookHeading,fill=BookPanel,rounded corners=1pt,text width=38mm,inner sep=5pt,align=left,font=\footnotesize},edge/.style={-{Latex[length=2mm]},draw=BookHeading}]
\node[box] (FORM_FLOW) at (0.0,0) {FORM\_FLOW};
\node[box] (FIELD) at (-2.5,-4) {\textbf{FIELD}\\\rule{\linewidth}{.3pt}\\string name\\string title\\string kind};
\node[box] (FIELD_RESULT) at (-2.5,-8) {FIELD\_RESULT};
\node[box] (GROUP_RESULT) at (2.5,-4) {\textbf{GROUP\_RESULT}\\\rule{\linewidth}{.3pt}\\boolean cancelled\\object values};
\node[box] (VALUE) at (2.5,-8) {VALUE};
\draw[edge,] (FORM_FLOW) to node[midway,fill=white,font=\scriptsize,text=BookMuted,align=center,text width=27mm] {contains / one : zero or more} (FIELD);
\draw[edge,] (FIELD) to node[midway,fill=white,font=\scriptsize,text=BookMuted,align=center,text width=27mm] {returns / one : one} (FIELD_RESULT);
\draw[edge,] (FORM_FLOW) to node[midway,fill=white,font=\scriptsize,text=BookMuted,align=center,text width=27mm] {aggregates / one : one} (GROUP_RESULT);
\draw[edge,] (GROUP_RESULT) to node[midway,fill=white,font=\scriptsize,text=BookMuted,align=center,text width=27mm] {includes / one : zero or more} (VALUE);
\end{tikzpicture}}
\end{center}

\section{Scenario 3: Make A Basic TUI App}\label{record-015-scenario-3-make-a-basic-tui-app}

Use the TEA runtime when the app needs a retained model, input handling, commands, animation, or repeated frames.

The loop is:

\begin{enumerate}
\def\labelenumi{\arabic{enumi}.}
\tightlist
\item
  \texttt{i\allowbreak{}n\allowbreak{}i\allowbreak{}t\allowbreak{}(\allowbreak{})} creates the first model and commands.
\item
  Runtime delivers messages.
\item
  \texttt{u\allowbreak{}p\allowbreak{}d\allowbreak{}a\allowbreak{}t\allowbreak{}e\allowbreak{}(\allowbreak{})} returns a new model and commands.
\item
  \texttt{v\allowbreak{}i\allowbreak{}e\allowbreak{}w\allowbreak{}(\allowbreak{})} lowers the model to a \texttt{S\allowbreak{}u\allowbreak{}r\allowbreak{}f\allowbreak{}a\allowbreak{}c\allowbreak{}e}.
\item
  The host writes the minimal terminal diff.
\end{enumerate}

\begin{center}
\resizebox{\linewidth}{!}{\begin{tikzpicture}[box/.style={draw=BookHeading,fill=BookPanel,rounded corners=1pt,text width=38mm,inner sep=5pt,align=left,font=\footnotesize},edge/.style={-{Latex[length=2mm]},draw=BookHeading}]
\node[box] (Init) at (0.0,0.0) {init};
\node[box] (Model) at (-2.5,-2.6) {Model};
\node[box] (Cmds) at (2.5,-2.6) {Commands};
\node[box] (Runtime) at (2.5,-10.4) {Runtime};
\node[box] (Msg) at (2.5,-13.0) {Message};
\node[box] (Update) at (0.0,-5.2) {update(msg, model)};
\node[box] (NewModel) at (-2.5,-7.800000000000001) {New model};
\node[box] (NewCmds) at (2.5,-7.800000000000001) {New commands};
\node[box] (View) at (-2.5,-10.4) {view(model)};
\node[box] (Surface) at (-2.5,-13.0) {Surface};
\node[box] (Diff) at (0.0,-15.600000000000001) {ANSI diff};
\node[box] (Terminal) at (0.0,-18.2) {Terminal};
\draw[edge,] (Init) to (Model);
\draw[edge,] (Init) to (Cmds);
\draw[edge,] (Runtime) to (Msg);
\draw[edge,] (Msg) to (Update);
\draw[edge,] (Model) to (Update);
\draw[edge,] (Update) to (NewModel);
\draw[edge,] (Update) to (NewCmds);
\draw[edge,] (NewModel) to (View);
\draw[edge,] (View) to (Surface);
\draw[edge,] (Surface) to (Diff);
\draw[edge,] (Diff) to (Terminal);
\draw[edge,] (NewCmds) to (Runtime);
\end{tikzpicture}}
\end{center}

\begin{center}
\resizebox{\linewidth}{!}{\begin{tikzpicture}[actor/.style={draw=BookHeading,fill=BookPanel,text width=27mm,align=center,font=\small,inner sep=5pt}]
\node[actor] (Terminal) at (0.0,0) {Terminal};
\draw[dashed,BookMuted] (Terminal.south) -- (0.0,-19.3);
\node[actor] (Runtime) at (4.5,0) {Runtime};
\draw[dashed,BookMuted] (Runtime.south) -- (4.5,-19.3);
\node[actor] (App) at (9.0,0) {App};
\draw[dashed,BookMuted] (App.south) -- (9.0,-19.3);
\draw[-{Latex[length=2mm]},BookHeading] (4.5,-1.7) -- node[above,fill=white,align=center,text width=4.05cm,font=\footnotesize] {init()} (9.0,-1.7);
\draw[dashed,-{Latex[length=2mm]},BookHeading] (9.0,-3.4) -- node[above,fill=white,align=center,text width=4.05cm,font=\footnotesize] {[model, commands]} (4.5,-3.4);
\draw[-{Latex[length=2mm]},BookHeading] (4.5,-5.1) -- node[above,fill=white,align=center,text width=4.05cm,font=\footnotesize] {view(model)} (9.0,-5.1);
\draw[dashed,-{Latex[length=2mm]},BookHeading] (9.0,-6.8) -- node[above,fill=white,align=center,text width=4.05cm,font=\footnotesize] {Surface} (4.5,-6.8);
\draw[-{Latex[length=2mm]},BookHeading] (4.5,-8.5) -- node[above,fill=white,align=center,text width=4.05cm,font=\footnotesize] {paint frame} (0.0,-8.5);
\draw[dashed,-{Latex[length=2mm]},BookHeading] (0.0,-10.2) -- node[above,fill=white,align=center,text width=4.05cm,font=\footnotesize] {key "+} (4.5,-10.2);
\draw[-{Latex[length=2mm]},BookHeading] (4.5,-11.9) -- node[above,fill=white,align=center,text width=4.05cm,font=\footnotesize] {update(key, model)} (9.0,-11.9);
\draw[dashed,-{Latex[length=2mm]},BookHeading] (9.0,-13.6) -- node[above,fill=white,align=center,text width=4.05cm,font=\footnotesize] {[new\allowbreak{}Model, commands]} (4.5,-13.6);
\draw[-{Latex[length=2mm]},BookHeading] (4.5,-15.299999999999999) -- node[above,fill=white,align=center,text width=4.05cm,font=\footnotesize] {view(new\allowbreak{}Model)} (9.0,-15.299999999999999);
\draw[dashed,-{Latex[length=2mm]},BookHeading] (9.0,-17.0) -- node[above,fill=white,align=center,text width=4.05cm,font=\footnotesize] {Surface} (4.5,-17.0);
\draw[-{Latex[length=2mm]},BookHeading] (4.5,-18.7) -- node[above,fill=white,align=center,text width=4.05cm,font=\footnotesize] {paint diff} (0.0,-18.7);
\end{tikzpicture}}\end{center}

\subsection{Code}\label{record-015-code-2}

\begin{CodeBlock}
import { box, kbd, stringToSurface } from '@flyingrobots/bijou';
import { startApp } from '@flyingrobots/bijou-node';
import { isKeyMsg, quit, type App } from '@flyingrobots/bijou-tui';

interface Model {
  readonly count: number;
}

type Msg = { type: 'increment' } | { type: 'decrement' } | { type: 'quit' };

const app: App<Model, Msg> = {
  init: () => [{ count: 0 }, []],

  update: (msg, model) => {
    if (isKeyMsg(msg)) {
      if (msg.key === '+' || msg.key === 'up') {
        return [{ count: model.count + 1 }, []];
      }
      if (msg.key === '-' || msg.key === 'down') {
        return [{ count: model.count - 1 }, []];
      }
      if (msg.key === 'q' || (msg.ctrl && msg.key === 'c')) {
        return [model, [quit()]];
      }
    }

    return [model, []];
  },

  view: (model) => {
    const text = [
      box(`Count: ${model.count}`),
      '',
      `${kbd('+')} increment`,
      `${kbd('-')} decrement`,
      `${kbd('q')} quit`,
    ].join('\n');

    return stringToSurface(text, 32, 8);
  },
};

await startApp(app);
\end{CodeBlock}

\subsection{Ownership Boundaries}\label{record-015-ownership-boundaries}

\begin{center}
\resizebox{\linewidth}{!}{\begin{tikzpicture}[box/.style={draw=BookHeading,fill=BookPanel,rounded corners=1pt,text width=38mm,inner sep=5pt,align=left,font=\footnotesize},edge/.style={-{Latex[length=2mm]},draw=BookHeading}]
\node[box] (Model) at (-5.0,0) {\textbf{Model}\\\rule{\linewidth}{.3pt}\\+count};
\node[box] (Msg) at (0.0,0) {\textbf{Msg}\\\rule{\linewidth}{.3pt}\\+key\\+type};
\node[box] (Update) at (5.0,0) {\textbf{Update}\\\rule{\linewidth}{.3pt}\\+update(msg, model)};
\node[box] (View) at (-2.5,-4) {\textbf{View}\\\rule{\linewidth}{.3pt}\\+view(model)};
\node[box] (Surface) at (2.5,-4) {\textbf{Surface}\\\rule{\linewidth}{.3pt}\\+cells};
\draw[edge,bend left=15] (Msg) to (Update);
\draw[edge,bend left=15] (Model) to (Update);
\draw[edge,bend left=15] (Update) to (Model);
\draw[edge,] (Model) to (View);
\draw[edge,bend left=15] (View) to (Surface);
\end{tikzpicture}}
\end{center}

Do not mutate the model from inside \texttt{v\allowbreak{}i\allowbreak{}e\allowbreak{}w\allowbreak{}(\allowbreak{})}. Views render. Updates own state change. Commands express runtime effects such as quit or delayed messages.

\section{Scenario 4: Make A Braille Ray Trace Renderer}\label{record-015-scenario-4-make-a-braille-ray-trace-renderer}

Use the TUI \texttt{c\allowbreak{}a\allowbreak{}n\allowbreak{}v\allowbreak{}a\allowbreak{}s\allowbreak{}(\allowbreak{})} renderer when you want procedural character graphics. The Braille resolution samples each terminal cell as a 2 by 4 sub-pixel grid, so a small terminal surface can carry more visual detail than one character per cell.

\begin{center}
\resizebox{\linewidth}{!}{\begin{tikzpicture}[box/.style={draw=BookHeading,fill=BookPanel,rounded corners=1pt,text width=38mm,inner sep=5pt,align=left,font=\footnotesize},edge/.style={-{Latex[length=2mm]},draw=BookHeading}]
\node[box] (Model) at (0.0,0.0) {Model: time and camera};
\node[box] (View) at (0.0,-2.6) {View};
\node[box] (Canvas) at (0.0,-5.2) {canvas(cols, rows, shader, resolution=braille)};
\node[box] (Samples) at (0.0,-7.800000000000001) {2x4 sub-pixel samples per cell};
\node[box] (Shader) at (0.0,-10.4) {Ray shader};
\node[box] (Cell) at (0.0,-13.0) {Cell color and on/off coverage};
\node[box] (Surface) at (0.0,-15.600000000000001) {Surface};
\node[box] (Terminal) at (0.0,-18.2) {Terminal};
\draw[edge,] (Model) to (View);
\draw[edge,] (View) to (Canvas);
\draw[edge,] (Canvas) to (Samples);
\draw[edge,] (Samples) to (Shader);
\draw[edge,] (Shader) to (Cell);
\draw[edge,] (Cell) to (Surface);
\draw[edge,] (Surface) to (Terminal);
\end{tikzpicture}}
\end{center}

\begin{center}
\resizebox{\linewidth}{!}{\begin{tikzpicture}[actor/.style={draw=BookHeading,fill=BookPanel,text width=27mm,align=center,font=\small,inner sep=5pt}]
\node[actor] (Runtime) at (0.0,0) {Runtime};
\draw[dashed,BookMuted] (Runtime.south) -- (0.0,-14.2);
\node[actor] (App) at (4.5,0) {App};
\draw[dashed,BookMuted] (App.south) -- (4.5,-14.2);
\node[actor] (Canvas) at (9.0,0) {Canvas};
\draw[dashed,BookMuted] (Canvas.south) -- (9.0,-14.2);
\node[actor] (Shader) at (13.5,0) {Shader};
\draw[dashed,BookMuted] (Shader.south) -- (13.5,-14.2);
\draw[-{Latex[length=2mm]},BookHeading] (0.0,-1.7) -- node[above,fill=white,align=center,text width=4.05cm,font=\footnotesize] {tick message} (4.5,-1.7);
\draw[dashed,-{Latex[length=2mm]},BookHeading] (4.5,-3.4) -- node[above,fill=white,align=center,text width=4.05cm,font=\footnotesize] {model with next time} (0.0,-3.4);
\draw[-{Latex[length=2mm]},BookHeading] (0.0,-5.1) -- node[above,fill=white,align=center,text width=4.05cm,font=\footnotesize] {view(model)} (4.5,-5.1);
\draw[-{Latex[length=2mm]},BookHeading] (4.5,-6.8) -- node[above,fill=white,align=center,text width=4.05cm,font=\footnotesize] {canvas(width, height, shader, braille)} (9.0,-6.8);
\draw[-{Latex[length=2mm]},BookHeading] (9.0,-8.5) -- node[above,fill=white,align=center,text width=4.05cm,font=\footnotesize] {\{ u, v, time, uniforms \}} (13.5,-8.5);
\draw[dashed,-{Latex[length=2mm]},BookHeading] (13.5,-10.2) -- node[above,fill=white,align=center,text width=4.05cm,font=\footnotesize] {cell or string} (9.0,-10.2);
\draw[dashed,-{Latex[length=2mm]},BookHeading] (9.0,-11.9) -- node[above,fill=white,align=center,text width=4.05cm,font=\footnotesize] {Surface} (4.5,-11.9);
\draw[dashed,-{Latex[length=2mm]},BookHeading] (4.5,-13.6) -- node[above,fill=white,align=center,text width=4.05cm,font=\footnotesize] {Surface} (0.0,-13.6);
\draw[dotted,BookHeading] (-.4,-7.6) rectangle (13.9,-10.5);
\node[anchor=west,font=\scriptsize,fill=white] at (-.3,-7.6) {Loop: each sub-pixel};
\end{tikzpicture}}\end{center}

\begin{center}
\resizebox{\linewidth}{!}{\begin{tikzpicture}[box/.style={draw=BookHeading,fill=BookPanel,rounded corners=1pt,text width=38mm,inner sep=5pt,align=left,font=\footnotesize},edge/.style={-{Latex[length=2mm]},draw=BookHeading}]
\node[box] (ShaderParams) at (-5.0,0) {\textbf{Shader\allowbreak{}Params}\\\rule{\linewidth}{.3pt}\\+u\\+v\\+time\\+uniforms};
\node[box] (ShaderCell) at (0.0,0) {\textbf{Shader\allowbreak{}Cell}\\\rule{\linewidth}{.3pt}\\+char\\+fg\allowbreak{}RGB\\+bg\allowbreak{}RGB};
\node[box] (ShaderFn) at (5.0,0) {\textbf{Shader\allowbreak{}Fn}\\\rule{\linewidth}{.3pt}\\+call(params)};
\node[box] (CanvasOptions) at (0.0,-4) {\textbf{Canvas\allowbreak{}Options}\\\rule{\linewidth}{.3pt}\\+resolution\\+time\\+uniforms};
\draw[edge,bend left=15] (ShaderFn) to (ShaderParams);
\draw[edge,bend left=15] (ShaderFn) to (ShaderCell);
\draw[edge,] (CanvasOptions) to (ShaderFn);
\end{tikzpicture}}
\end{center}

\subsection{Code}\label{record-015-code-3}

\begin{CodeBlock}
import { startApp } from '@flyingrobots/bijou-node';
import {
  canvas,
  isKeyMsg,
  quit,
  tick,
  type App,
  type ShaderFn,
} from '@flyingrobots/bijou-tui';

interface Model {
  readonly time: number;
}

type Msg = { type: 'tick' };

function sphereDepth(x: number, y: number, radius: number): number | null {
  const d2 = x * x + y * y;
  if (d2 > radius * radius) return null;
  return Math.sqrt(radius * radius - d2);
}

const shader: ShaderFn = ({ u, v, time }) => {
  const x = (u - 0.5) * 2.4;
  const y = (v - 0.5) * 1.8;
  const z = sphereDepth(x, y, 0.82);

  if (z === null) {
    return { char: ' ', bgRGB: [8, 10, 22] };
  }

  const lightX = Math.cos(time * 0.8) * 0.55;
  const lightY = -0.35;
  const lightZ = 0.9;
  const lightLength = Math.hypot(lightX, lightY, lightZ);
  const shade = Math.max(0, (x * lightX + y * lightY + z * lightZ) / lightLength);
  const rim = Math.max(0, 1 - z);

  return {
    char: shade + rim * 0.2 > 0.28 ? 'x' : ' ',
    fgRGB: [
      Math.round(80 + shade * 150),
      Math.round(110 + shade * 120),
      Math.round(180 + shade * 60),
    ],
    bgRGB: [8, 10, 22],
  };
};

const app: App<Model, Msg> = {
  init: () => [{ time: 0 }, [tick(33, { type: 'tick' })]],

  update: (msg, model) => {
    if (isKeyMsg(msg) && (msg.key === 'q' || (msg.ctrl && msg.key === 'c'))) {
      return [model, [quit()]];
    }

    if ('type' in msg && msg.type === 'tick') {
      return [{ time: model.time + 0.033 }, [tick(33, { type: 'tick' })]];
    }

    return [model, []];
  },

  view: (model) =>
    canvas(72, 28, shader, {
      resolution: 'braille',
      time: model.time,
    }),
};

await startApp(app);
\end{CodeBlock}

\subsection{Renderer Data Flow}\label{record-015-renderer-data-flow}

\begin{center}
\resizebox{\linewidth}{!}{\begin{tikzpicture}[box/.style={draw=BookHeading,fill=BookPanel,rounded corners=1pt,text width=38mm,inner sep=5pt,align=left,font=\footnotesize},edge/.style={-{Latex[length=2mm]},draw=BookHeading}]
\node[box] (MODEL) at (-2.5,0.0) {MODEL};
\node[box] (VIEW) at (-2.5,-2.6) {VIEW};
\node[box] (CANVAS_CALL) at (0.0,-5.2) {CANVAS\_CALL};
\node[box] (SHADER_SAMPLE) at (0.0,-7.800000000000001) {SHADER\_SAMPLE};
\node[box] (SHADER_CELL) at (0.0,-10.4) {SHADER\_CELL};
\node[box] (SURFACE_CELL) at (2.5,-2.6) {SURFACE\_CELL};
\node[box] (SURFACE) at (2.5,0.0) {SURFACE};
\draw[edge,] (MODEL) to node[midway,fill=white,font=\scriptsize,text=BookMuted,align=center,text width=27mm] {drives / one : one} (VIEW);
\draw[edge,] (VIEW) to node[midway,fill=white,font=\scriptsize,text=BookMuted,align=center,text width=27mm] {creates / one : one} (CANVAS_CALL);
\draw[edge,] (CANVAS_CALL) to node[midway,fill=white,font=\scriptsize,text=BookMuted,align=center,text width=27mm] {emits / one : zero or more} (SHADER_SAMPLE);
\draw[edge,] (SHADER_SAMPLE) to node[midway,fill=white,font=\scriptsize,text=BookMuted,align=center,text width=27mm] {returns / one : one} (SHADER_CELL);
\draw[edge,] (SHADER_CELL) to node[midway,fill=white,font=\scriptsize,text=BookMuted,align=center,text width=27mm] {accumulates\_into / zero or more : one} (SURFACE_CELL);
\draw[edge,] (SURFACE) to node[midway,fill=white,font=\scriptsize,text=BookMuted,align=center,text width=27mm] {contains / one : zero or more} (SURFACE_CELL);
\end{tikzpicture}}
\end{center}

\subsection{Performance Notes}\label{record-015-performance-notes}

\begin{itemize}
\tightlist
\item
  Keep shader math deterministic and bounded.
\item
  Keep allocations out of the inner shader path.
\item
  Use \texttt{u\allowbreak{}n\allowbreak{}i\allowbreak{}f\allowbreak{}o\allowbreak{}r\allowbreak{}m\allowbreak{}s} for static scene values instead of closing over mutable runtime objects.
\item
  Prefer \texttt{r\allowbreak{}e\allowbreak{}s\allowbreak{}o\allowbreak{}l\allowbreak{}u\allowbreak{}t\allowbreak{}i\allowbreak{}o\allowbreak{}n\allowbreak{}:\allowbreak{} \allowbreak{}'\allowbreak{}b\allowbreak{}r\allowbreak{}a\allowbreak{}i\allowbreak{}l\allowbreak{}l\allowbreak{}e\allowbreak{}'} for high-density binary coverage and \texttt{r\allowbreak{}e\allowbreak{}s\allowbreak{}o\allowbreak{}l\allowbreak{}u\allowbreak{}t\allowbreak{}i\allowbreak{}o\allowbreak{}n\allowbreak{}:\allowbreak{} \allowbreak{}'\allowbreak{}g\allowbreak{}l\allowbreak{}y\allowbreak{}p\allowbreak{}h\allowbreak{}'} when you want glyph fitting from coverage.
\item
  Use \texttt{n\allowbreak{}p\allowbreak{}m\allowbreak{} \allowbreak{}r\allowbreak{}u\allowbreak{}n\allowbreak{} \allowbreak{}b\allowbreak{}e\allowbreak{}n\allowbreak{}c\allowbreak{}h} and \texttt{n\allowbreak{}p\allowbreak{}m\allowbreak{} \allowbreak{}r\allowbreak{}u\allowbreak{}n\allowbreak{} \allowbreak{}s\allowbreak{}o\allowbreak{}a\allowbreak{}k} when changing renderer internals.
\end{itemize}

\section{Mode Lowering}\label{record-015-mode-lowering}

Bijou components should degrade intentionally when output is not an interactive terminal.

\begin{center}
\resizebox{\linewidth}{!}{\begin{tikzpicture}[box/.style={draw=BookHeading,fill=BookPanel,rounded corners=1pt,text width=38mm,inner sep=5pt,align=left,font=\footnotesize},edge/.style={-{Latex[length=2mm]},draw=BookHeading}]
\node[box] (Component) at (0.0,0.0) {Component or Block};
\node[box] (Interactive) at (-7.5,-2.6) {interactive: rich terminal UI};
\node[box] (Static) at (-2.5,-2.6) {static: fixed visual output};
\node[box] (Pipe) at (2.5,-2.6) {pipe: plain text};
\node[box] (Accessible) at (7.5,-2.6) {accessible: screen-reader text};
\node[box] (Facts) at (0.0,-5.2) {semantic facts};
\draw[edge,] (Component) to (Interactive);
\draw[edge,] (Component) to (Static);
\draw[edge,] (Component) to (Pipe);
\draw[edge,] (Component) to (Accessible);
\draw[edge,] (Interactive) to (Facts);
\draw[edge,] (Static) to (Facts);
\draw[edge,] (Pipe) to (Facts);
\draw[edge,] (Accessible) to (Facts);
\end{tikzpicture}}
\end{center}

Use this rule:

\begin{itemize}
\tightlist
\item
  Interactive and static modes may draw boxes, colors, charts, or Braille surfaces.
\item
  Pipe mode should emit useful plain text.
\item
  Accessible mode should prioritize meaningful spoken order.
\item
  Semantic facts should survive lowering when tooling or agents need to inspect the output.
\end{itemize}

\section{Blocks And Data Binding}\label{record-015-blocks-and-data-binding}

For reusable application surfaces, use Blocks. Blocks declare metadata, variants, slots, data requirements, command intents, and mode-lowering posture. They do not pull mutable state directly.

\begin{center}
\resizebox{\linewidth}{!}{\begin{tikzpicture}[box/.style={draw=BookHeading,fill=BookPanel,rounded corners=1pt,text width=38mm,inner sep=5pt,align=left,font=\footnotesize},edge/.style={-{Latex[length=2mm]},draw=BookHeading}]
\node[box] (Provider) at (0.0,0.0) {Provider publishes immutable snapshot};
\node[box] (Snapshot) at (0.0,-2.6) {Binding\allowbreak{}Snapshot};
\node[box] (Frame) at (0.0,-5.2) {Binding\allowbreak{}Frame};
\node[box] (Block) at (0.0,-7.800000000000001) {Block render};
\node[box] (Surface) at (-2.5,-10.4) {Surface or lowered output};
\node[box] (Intent) at (2.5,-10.4) {Command\allowbreak{}Intent};
\node[box] (Business) at (0.0,-13.0) {Business logic};
\draw[edge,] (Provider) to (Snapshot);
\draw[edge,] (Snapshot) to (Frame);
\draw[edge,] (Frame) to (Block);
\draw[edge,] (Block) to (Surface);
\draw[edge,] (Block) to (Intent);
\draw[edge,] (Intent) to (Business);
\draw[edge,] (Business) to (Provider);
\end{tikzpicture}}
\end{center}

\begin{center}
\resizebox{\linewidth}{!}{\begin{tikzpicture}[box/.style={draw=BookHeading,fill=BookPanel,rounded corners=1pt,text width=38mm,inner sep=5pt,align=left,font=\footnotesize},edge/.style={-{Latex[length=2mm]},draw=BookHeading}]
\node[box] (DataRequirement) at (-5.0,0) {\textbf{Data\allowbreak{}Requirement}\\\rule{\linewidth}{.3pt}\\+id\\+resource};
\node[box] (BindingSnapshot) at (0.0,0) {\textbf{Binding\allowbreak{}Snapshot}\\\rule{\linewidth}{.3pt}\\+provider\allowbreak{}Id\\+requirement\allowbreak{}Id\\+version\\+status\\+data};
\node[box] (BindingFrame) at (5.0,0) {\textbf{Binding\allowbreak{}Frame}\\\rule{\linewidth}{.3pt}\\+require(id)\\+get(id)\\+status(id)\\+issues(id)};
\node[box] (CommandIntent) at (-2.5,-4) {\textbf{Command\allowbreak{}Intent}\\\rule{\linewidth}{.3pt}\\+id\\+name};
\node[box] (BlockDefinition) at (2.5,-4) {\textbf{Block\allowbreak{}Definition}\\\rule{\linewidth}{.3pt}\\+metadata\\+render(input)};
\draw[edge,bend left=15] (DataRequirement) to (BindingSnapshot);
\draw[edge,bend left=15] (BindingSnapshot) to (BindingFrame);
\draw[edge,] (BindingFrame) to (BlockDefinition);
\draw[edge,bend left=15] (BlockDefinition) to (CommandIntent);
\end{tikzpicture}}
\end{center}

That loop is unidirectional:

\begin{CodeBlock}
business logic/providers
  -> immutable snapshots
  -> render frame
  -> views/blocks render
  -> command intents leave as user intent
  -> business logic owns change
\end{CodeBlock}

\section{Common Mistakes}\label{record-015-common-mistakes}

\begin{center}
\resizebox{\linewidth}{!}{\begin{tikzpicture}[box/.style={draw=BookHeading,fill=BookPanel,rounded corners=1pt,text width=38mm,inner sep=5pt,align=left,font=\footnotesize},edge/.style={-{Latex[length=2mm]},draw=BookHeading}]
\node[box] (Bad1) at (-10.0,0.0) {Render reads files or process.env};
\node[box] (Fix1) at (-10.0,-2.6) {Move host access behind adapter};
\node[box] (Bad2) at (-5.0,0.0) {View mutates model};
\node[box] (Fix2) at (-5.0,-2.6) {Return new model from update};
\node[box] (Bad3) at (0.0,0.0) {Block pulls provider directly};
\node[box] (Fix3) at (0.0,-2.6) {Use snapshots and Binding\allowbreak{}Frame};
\node[box] (Bad4) at (5.0,0.0) {Canvas allocates per sample};
\node[box] (Fix4) at (5.0,-2.6) {Precompute uniforms};
\node[box] (Bad5) at (10.0,0.0) {CLI starts full-screen runtime for one line};
\node[box] (Fix5) at (10.0,-2.6) {Use core components};
\draw[edge,] (Bad1) to (Fix1);
\draw[edge,] (Bad2) to (Fix2);
\draw[edge,] (Bad3) to (Fix3);
\draw[edge,] (Bad4) to (Fix4);
\draw[edge,] (Bad5) to (Fix5);
\end{tikzpicture}}
\end{center}

{\def\LTcaptype{none} % do not increment counter
\begin{longtable}[]{@{}
  >{\raggedright\arraybackslash}p{(\linewidth - 2\tabcolsep) * \real{0.5000}}
  >{\raggedright\arraybackslash}p{(\linewidth - 2\tabcolsep) * \real{0.5000}}@{}}
\toprule\noalign{}
{M\allowbreak{}i\allowbreak{}s\allowbreak{}t\allowbreak{}a\allowbreak{}k\allowbreak{}e}
 & {B\allowbreak{}e\allowbreak{}t\allowbreak{}t\allowbreak{}e\allowbreak{}r} {p\allowbreak{}a\allowbreak{}t\allowbreak{}h}
 \\
\midrule\noalign{}
\endhead
\bottomrule\noalign{}
\endlastfoot
{P\allowbreak{}a\allowbreak{}r\allowbreak{}s\allowbreak{}i\allowbreak{}n\allowbreak{}g} \texttt{s\allowbreak{}t\allowbreak{}d\allowbreak{}i\allowbreak{}n} {i\allowbreak{}n\allowbreak{}s\allowbreak{}i\allowbreak{}d\allowbreak{}e} {a} {c\allowbreak{}o\allowbreak{}m\allowbreak{}p\allowbreak{}o\allowbreak{}n\allowbreak{}e\allowbreak{}n\allowbreak{}t} & {L\allowbreak{}e\allowbreak{}t} {h\allowbreak{}o\allowbreak{}s\allowbreak{}t\allowbreak{}/\allowbreak{}r\allowbreak{}u\allowbreak{}n\allowbreak{}t\allowbreak{}i\allowbreak{}m\allowbreak{}e} {p\allowbreak{}r\allowbreak{}o\allowbreak{}d\allowbreak{}u\allowbreak{}c\allowbreak{}e} {m\allowbreak{}e\allowbreak{}s\allowbreak{}s\allowbreak{}a\allowbreak{}g\allowbreak{}e\allowbreak{}s} {o\allowbreak{}r} {f\allowbreak{}o\allowbreak{}r\allowbreak{}m} {a\allowbreak{}n\allowbreak{}s\allowbreak{}w\allowbreak{}e\allowbreak{}r\allowbreak{}s\allowbreak{}.} \\
{C\allowbreak{}a\allowbreak{}l\allowbreak{}l\allowbreak{}i\allowbreak{}n\allowbreak{}g} {p\allowbreak{}r\allowbreak{}o\allowbreak{}v\allowbreak{}i\allowbreak{}d\allowbreak{}e\allowbreak{}r} {r\allowbreak{}e\allowbreak{}f\allowbreak{}r\allowbreak{}e\allowbreak{}s\allowbreak{}h} {f\allowbreak{}r\allowbreak{}o\allowbreak{}m} {r\allowbreak{}e\allowbreak{}n\allowbreak{}d\allowbreak{}e\allowbreak{}r} & {E\allowbreak{}m\allowbreak{}i\allowbreak{}t} {a} {c\allowbreak{}o\allowbreak{}m\allowbreak{}m\allowbreak{}a\allowbreak{}n\allowbreak{}d} {i\allowbreak{}n\allowbreak{}t\allowbreak{}e\allowbreak{}n\allowbreak{}t\allowbreak{};} {b\allowbreak{}u\allowbreak{}s\allowbreak{}i\allowbreak{}n\allowbreak{}e\allowbreak{}s\allowbreak{}s} {l\allowbreak{}o\allowbreak{}g\allowbreak{}i\allowbreak{}c} {o\allowbreak{}w\allowbreak{}n\allowbreak{}s} {c\allowbreak{}h\allowbreak{}a\allowbreak{}n\allowbreak{}g\allowbreak{}e\allowbreak{}.} \\
{D\allowbreak{}u\allowbreak{}p\allowbreak{}l\allowbreak{}i\allowbreak{}c\allowbreak{}a\allowbreak{}t\allowbreak{}i\allowbreak{}n\allowbreak{}g} {E\allowbreak{}n\allowbreak{}g\allowbreak{}l\allowbreak{}i\allowbreak{}s\allowbreak{}h} {f\allowbreak{}a\allowbreak{}l\allowbreak{}l\allowbreak{}b\allowbreak{}a\allowbreak{}c\allowbreak{}k} {s\allowbreak{}t\allowbreak{}r\allowbreak{}i\allowbreak{}n\allowbreak{}g\allowbreak{}s} {i\allowbreak{}n} {t\allowbreak{}r\allowbreak{}a\allowbreak{}n\allowbreak{}s\allowbreak{}l\allowbreak{}a\allowbreak{}t\allowbreak{}e\allowbreak{}d} {c\allowbreak{}a\allowbreak{}t\allowbreak{}a\allowbreak{}l\allowbreak{}o\allowbreak{}g\allowbreak{}s} & {L\allowbreak{}o\allowbreak{}a\allowbreak{}d} {E\allowbreak{}n\allowbreak{}g\allowbreak{}l\allowbreak{}i\allowbreak{}s\allowbreak{}h} {f\allowbreak{}a\allowbreak{}l\allowbreak{}l\allowbreak{}b\allowbreak{}a\allowbreak{}c\allowbreak{}k} {s\allowbreak{}e\allowbreak{}p\allowbreak{}a\allowbreak{}r\allowbreak{}a\allowbreak{}t\allowbreak{}e\allowbreak{}l\allowbreak{}y\allowbreak{}.} \\
{U\allowbreak{}s\allowbreak{}i\allowbreak{}n\allowbreak{}g} {a} {f\allowbreak{}u\allowbreak{}l\allowbreak{}l\allowbreak{}-\allowbreak{}s\allowbreak{}c\allowbreak{}r\allowbreak{}e\allowbreak{}e\allowbreak{}n} {T\allowbreak{}U\allowbreak{}I} {f\allowbreak{}o\allowbreak{}r} {s\allowbreak{}i\allowbreak{}m\allowbreak{}p\allowbreak{}l\allowbreak{}e} {o\allowbreak{}u\allowbreak{}t\allowbreak{}p\allowbreak{}u\allowbreak{}t} & {U\allowbreak{}s\allowbreak{}e} {c\allowbreak{}o\allowbreak{}r\allowbreak{}e} {c\allowbreak{}o\allowbreak{}m\allowbreak{}p\allowbreak{}o\allowbreak{}n\allowbreak{}e\allowbreak{}n\allowbreak{}t\allowbreak{}s} {a\allowbreak{}n\allowbreak{}d} {a} {N\allowbreak{}o\allowbreak{}d\allowbreak{}e} {c\allowbreak{}o\allowbreak{}n\allowbreak{}t\allowbreak{}e\allowbreak{}x\allowbreak{}t\allowbreak{}.} \\
{T\allowbreak{}r\allowbreak{}e\allowbreak{}a\allowbreak{}t\allowbreak{}i\allowbreak{}n\allowbreak{}g} {t\allowbreak{}e\allowbreak{}r\allowbreak{}m\allowbreak{}i\allowbreak{}n\allowbreak{}a\allowbreak{}l} {g\allowbreak{}r\allowbreak{}a\allowbreak{}p\allowbreak{}h\allowbreak{}i\allowbreak{}c\allowbreak{}s} {a\allowbreak{}s} {r\allowbreak{}a\allowbreak{}w} {s\allowbreak{}t\allowbreak{}r\allowbreak{}i\allowbreak{}n\allowbreak{}g\allowbreak{}s} {e\allowbreak{}v\allowbreak{}e\allowbreak{}r\allowbreak{}y\allowbreak{}w\allowbreak{}h\allowbreak{}e\allowbreak{}r\allowbreak{}e} & {P\allowbreak{}r\allowbreak{}o\allowbreak{}m\allowbreak{}o\allowbreak{}t\allowbreak{}e} {t\allowbreak{}h\allowbreak{}e\allowbreak{}m} {t\allowbreak{}o} \texttt{S\allowbreak{}u\allowbreak{}r\allowbreak{}f\allowbreak{}a\allowbreak{}c\allowbreak{}e} {w\allowbreak{}h\allowbreak{}e\allowbreak{}n} {l\allowbreak{}a\allowbreak{}y\allowbreak{}o\allowbreak{}u\allowbreak{}t} {a\allowbreak{}n\allowbreak{}d} {d\allowbreak{}i\allowbreak{}f\allowbreak{}f\allowbreak{}i\allowbreak{}n\allowbreak{}g} {m\allowbreak{}a\allowbreak{}t\allowbreak{}t\allowbreak{}e\allowbreak{}r\allowbreak{}.} \\
\end{longtable}
}

\section{Validation Checklist}\label{record-015-validation-checklist}

Before shipping a terminal scenario, run the narrow command for the surface you changed, then the shared gates:

\begin{CodeBlock}
npm run lint
npm run typecheck:test
npm test
npm run docs:inventory
npm run verify:interactive-examples
\end{CodeBlock}

For user-facing product proof:

\begin{CodeBlock}
npm run dogfood
npm run blocklab
\end{CodeBlock}

For rendering or long-running runtime changes:

\begin{CodeBlock}
npm run bench
npm run soak
\end{CodeBlock}

